% ================================================================
% SECTION 6: CONSTRAINTS MODULE
% ================================================================
\section{Constraints Module}
\label{sec:constraints}

\begin{center}
\code{bionetflux.core.constraints} ~~and~~ \code{bionetflux.core.boundary\_override}
\end{center}

The constraints module (612 lines) manages Lagrange-multiplier-based enforcement of boundary and junction conditions.  The boundary override sub-module (187 lines) applies per-point, per-equation overrides from TOML configuration.

% ------------------------------------------------------------------
\subsection{ConstraintType Enum}

\begin{lstlisting}[language=Python, caption={Constraint types}]
class ConstraintType(Enum):
    DIRICHLET        = "dirichlet"
    NEUMANN          = "neumann"
    ROBIN            = "robin"
    TRACE_CONTINUITY = "trace_continuity"
    KEDEM_KATCHALSKY = "kedem_katchalsky"
\end{lstlisting}

The first three are \textbf{boundary conditions} (single domain, single multiplier); the last two are \textbf{junction conditions} (two domains, two multipliers).

% ------------------------------------------------------------------
\subsection{Constraint Class}

\begin{lstlisting}[language=Python, caption={Constraint constructor}]
class Constraint:
    def __init__(self, 
                 constraint_type: ConstraintType,
                 equation_index: int,
                 domains: List[int],
                 positions: List[float],
                 parameters: Optional[np.ndarray] = None,
                 data_function: Optional[Callable] = None):
\end{lstlisting}

\noindent Validation at construction time:
\begin{itemize}
  \item \code{len(domains)} must equal \code{len(positions)}.
  \item Boundary types require exactly one domain; junction types require exactly two.
\end{itemize}

\noindent Properties and methods:
\begin{longtable}{p{6cm}p{9cm}}
\toprule
\textbf{Member} & \textbf{Description} \\
\midrule
\code{is\_boundary\_condition} & \code{True} for Dirichlet, Neumann, Robin \\
\code{is\_junction\_condition} & \code{True} for Trace Continuity, Kedem-Katchalsky \\
\code{n\_multipliers -> int} & 1 for boundary, 2 for junction \\
\code{get\_data(time) -> float} & Evaluate the \code{data\_function} at time $t$ \\
\bottomrule
\end{longtable}

% ------------------------------------------------------------------
\subsection{ConstraintManager Class}

\begin{lstlisting}[language=Python, caption={ConstraintManager constructor}]
class ConstraintManager:
    def __init__(self):
        self.constraints: List[Constraint] = []
        self._node_mappings = []           # filled by map_to_discretizations
\end{lstlisting}

\subsubsection{Adding Constraints}

Two parallel API styles are provided: \code{add\_*} methods append and return an index; \code{make\_*} methods return a \code{Constraint} object without appending.

\begin{longtable}{p{7.5cm}p{7cm}}
\toprule
\textbf{Method} & \textbf{Description} \\
\midrule
\code{add\_constraint(constraint) -> int} & Append raw \code{Constraint} \\
\code{add\_dirichlet(eq\_idx, domain\_id, position, data\_function=None)} & Dirichlet BC \\
\code{add\_neumann(eq\_idx, domain\_id, position, data\_function=None)} & Neumann BC \\
\code{add\_robin(eq\_idx, domain\_id, position, alpha, beta, data\_function=None)} & Robin BC ($\alpha u + \beta \sigma = g$) \\
\code{add\_trace\_continuity(eq\_idx, dom1, dom2, pos1, pos2)} & Trace jump = 0 \\
\code{add\_kedem\_katchalsky(eq\_idx, dom1, dom2, pos1, pos2, permeability)} & Kedem-Katchalsky \\
\midrule
\code{make\_dirichlet(...) -> Constraint} & Factory (no append) \\
\code{make\_neumann(...) -> Constraint} & Factory (no append) \\
\code{make\_robin(...) -> Constraint} & Factory (no append) \\
\code{make\_trace\_continuity(...) -> Constraint} & Factory (no append) \\
\code{make\_kedem\_katchalsky(...) -> Constraint} & Factory (no append) \\
\bottomrule
\end{longtable}

\subsubsection{Querying and Modifying}

\begin{longtable}{p{8cm}p{7cm}}
\toprule
\textbf{Method} & \textbf{Description} \\
\midrule
\code{find\_constraints(domain\_index=None, equation\_index=None, constraint\_type=None, position=None) -> List[int]} & Find indices matching all supplied filters \\
\code{replace\_constraint(index, new\_constraint)} & Replace constraint at given index \\
\code{get\_constraints\_by\_domain(domain\_index) -> List[int]} & Filter by domain \\
\code{get\_constraints\_by\_type(constraint\_type) -> List[int]} & Filter by type \\
\bottomrule
\end{longtable}

\subsubsection{Assembly Interface}

\begin{longtable}{p{8cm}p{7cm}}
\toprule
\textbf{Method} & \textbf{Description} \\
\midrule
\code{map\_to\_discretizations(discretizations)} & Compute node indices from positions; must be called after all constraints are added and before assembly. \\
\code{get\_node\_indices(constraint\_index) -> List[int]} & Node indices for a constraint \\
\code{n\_constraints -> int} & Number of constraints (property) \\
\code{n\_multipliers -> int} & Total multiplier count (property) \\
\code{get\_multiplier\_data(time) -> np.ndarray} & Evaluate all data functions at time $t$ \\
\code{compute\_constraint\_residuals(...)} & Compute residual and Jacobian contributions for all constraints \\
\bottomrule
\end{longtable}

% ------------------------------------------------------------------
\subsection{Boundary Override Mechanism}
\label{sec:boundary_override}

\begin{center}
\code{bionetflux.core.boundary\_override}
\end{center}

The function \code{apply\_boundary\_overrides} replaces default Neumann constraints with overrides specified in the TOML \code{[boundary\_conditions]} section.

\begin{lstlisting}[language=Python, caption={Boundary override signature}]
def apply_boundary_overrides(
    constraint_manager: ConstraintManager,
    boundary_overrides: Dict[str, Any],
    boundary_point_map: Dict[str, tuple],
    equation_names: List[str],
    function_resolver: Optional[Any] = None,
) -> None:
\end{lstlisting}

\noindent The TOML key format is \code{<point\_name>\_<equation\_name>} and the value is an inline table:

\begin{lstlisting}[caption={Boundary override TOML format}]
[boundary_conditions]
B1_u   = { type = "dirichlet", data = "zeros" }
B2_phi = { type = "neumann",   data = "sin_t" }
B3_v   = { type = "robin",  alpha = 1.0, beta = 0.5, data = "zeros" }
\end{lstlisting}

\noindent The function:
\begin{enumerate}
  \item Parses the key to extract the point name and equation name.
  \item Looks up \code{(domain\_id, position)} from \code{boundary\_point\_map} (stored in geometry global metadata by \code{create\_maze\_geometry}).
  \item Finds the existing default Neumann constraint via \code{find\_constraints}.
  \item Creates the replacement constraint using the \code{make\_*} factory.
  \item Calls \code{replace\_constraint} to swap it in.
\end{enumerate}

The caller must call \code{constraint\_manager.map\_to\_discretizations()} after overrides are applied.
